% =====================================================================
% APPENDIX B: USE OF AI CODE ASSISTANTS
% =====================================================================
\appendixbanner{APPENDIX B: USE OF AI CODE ASSISTANTS}
\addcontentsline{toc}{section}{Appendix B: Use of AI code assistants}

\subsection*{Workflow}

The code was co-developed with Claude (Anthropic) working on the project folder through the Claude
desktop app. The division of work was fixed from the start: the author defined the dataset, the tasks,
the models and the validation protocol (listed in \texttt{docs/DECISIONS.md} with their justification);
the assistant wrote the code stage by stage, executed every notebook in a clean copy of the repository
and inspected every figure before delivering it; the author ran each notebook on her own computer, and
the assistant compared the resulting JSON summaries with its own run, value by value, before the next
stage. Every exchange that shaped the code is recorded in \texttt{docs/AI\_USAGE\_LOG.md}.

\subsection*{Errors detected (candidates for the AI Error Bounty)}

All of them were produced by the assistant and caught by the verification loop before delivery.

\begin{enumerate}
  \item \textbf{Non-reproducible results.} \texttt{LogisticRegression(solver="liblinear")} was created
        without \texttt{random\_state}; \texttt{liblinear} shuffles the data, so two runs chose different
        hyperparameters in the outer folds. \emph{Detection}: comparing two runs. \emph{Fix}:
        \texttt{random\_state=42}; two complete runs then gave identical out-of-fold probabilities.
  \item \textbf{Silent NaN in the regression diagnostics.} The summary reported ``maximum Cook's
        distance = nan''. \emph{Cause}: one race category has a single patient, so its one-hot column
        gives that patient leverage 1, the model fits it exactly, and the standardised residual is 0/0.
        \emph{Fix}: detect leverage-1 patients, report them and exclude them from the diagnostics.
  \item \textbf{Distorted biplot.} The loading arrows were scaled with a different factor on each axis,
        which changes their directions. \emph{Detection}: arrows missing from the scaled biplot.
        \emph{Fix}: one common factor.
  \item \textbf{Misleading tree figures.} sklearn's default colours painted LGG in orange and GBM in
        blue, the opposite of every other figure; and with zero decimals the split ``IDH1 $\leq$ 0.5''
        was printed as ``IDH1 $\leq$ 0.0''. \emph{Detection}: comparing the figure with the text export of
        the tree. \emph{Fix}: recolouring and one decimal.
  \item \textbf{Axis hiding data.} The number of mutations per patient was drawn on a categorical axis
        that silently skipped empty categories, hiding the few patients with many mutations.
        \emph{Fix}: grouped ``6+'' category.
  \item \textbf{Escaping errors in generated tables} (\texttt{\$\textbackslash pm\$} printed literally)
        and minor bugs (a negative error-bar length from floating-point rounding, a lost column name,
        NumPy~2 number formatting) found when executing the notebooks.
\end{enumerate}

\subsection*{Reflection}

The assistant wrote code that ran without errors and looked plausible in every one of these cases; the
errors were only visible by running the code twice, reading the figures against the numbers, or
connecting a NaN with the theory of high-leverage points. Owning the code meant owning this
verification loop.
